\documentclass[../cours_1re.tex]{subfiles}
\begin{document}
\setcounter{chapter}{11}
\chapter{La fonction exponentielle}
\label{ch:12}
\begin{objectif}
Définir l'exponentielle par une équation différentielle et démontrer ses propriétés algébriques et analytiques. L'existence est d'abord admise ; une construction complète est donnée au chapitre 26.
\end{objectif}
\section{Une définition par un taux de variation}
\begin{definition}
La fonction exponentielle est l'unique fonction dérivable sur $\R$ telle que
$\exp'(x)=\exp(x)$ pour tout $x$ et $\exp(0)=1$.
\end{definition}
Son accroissement instantané est proportionnel à sa propre valeur, avec coefficient $1$. On admet ici l'existence d'une telle fonction ; l'unicité et les propriétés seront démontrées. Le chapitre 26 construira l'objet, afin de fermer cette dépendance.
\section{L'exponentielle ne s'annule pas}
\begin{demonstration}
Poser $g(x)=\exp(x)\exp(-x)$. La règle du produit et la composition affine donnent
$g'(x)=\exp(x)\exp(-x)-\exp(x)\exp(-x)=0$.
Comme $\R$ est un intervalle, $g$ est constante ; $g(0)=1$, donc
\[
\exp(x)\exp(-x)=1.
\]
Aucun facteur ne peut donc être nul.
\end{demonstration}
\section{Unicité et addition des arguments}
\begin{demonstration}[Unicité]
Si $f'=f$ et $f(0)=1$, le quotient $f/\exp$ est défini partout et sa dérivée est
$(f'\exp-f\exp')/\exp^2=0$. Il est constant et vaut $1$ en zéro ; donc $f=\exp$.
\end{demonstration}
\begin{theoreme}
Pour tous réels $x,y$,
$\exp(x+y)=\exp(x)\exp(y)$.
\end{theoreme}
\begin{demonstration}
Fixer $y$ et considérer $h(x)=\exp(x+y)/\exp(x)$.
La règle du quotient donne $h'(x)=0$. Ainsi $h(x)=h(0)=\exp(y)$, d'où l'identité.
Le dénominateur est non nul par la preuve précédente.
\end{demonstration}
En particulier $\exp(x)=\exp(x/2)^2>0$. Puisque $\exp'=\exp>0$, la fonction est strictement croissante.
La preuve n'a pas supposé la positivité pour justifier une division : la non-annulation avait été établie séparément.
\section{Notation et règles de calcul}
On pose $e=\exp(1)$ et on écrit $e^x=\exp(x)$. Pour les entiers, cette notation coïncide avec les puissances usuelles, car l'addition répétée des arguments donne $\exp(n)=e^n$.
Pour $q=m/n$ rationnel, $n>0$, $\exp(q)^n=\exp(m)=e^m$, ce qui prolonge les puissances par racines positives.
Pour un exposant réel quelconque, la notation désigne la fonction construite, pas une multiplication répétée un nombre irrationnel de fois.
\[
e^{x+y}=e^xe^y,\quad e^{-x}=1/e^x,\quad
e^{x-y}=e^x/e^y,\quad (e^x)^n=e^{nx}\quad(n\in\Z).
\]
\begin{exemple}
$e^{2x}e^{3-x}/e^x=e^3$.
$e^{2x-1}=e^{x+4}$ équivaut à $2x-1=x+4$, donc $x=5$.
$e^{2x-1}<1$ équivaut à $2x-1<0$, donc $x<1/2$.
En revanche $e^{x+y}$ n'est pas $e^x+e^y$ : en $x=y=0$, les valeurs seraient $1$ et $2$.
\end{exemple}
\section{Courbe, tangente et croissance}
La courbe passe par $(0;1)$, avec une tangente $y=1+x$. La fonction
$g(x)=e^x-1-x$ a pour dérivée $e^x-1$, négative pour $x<0$ et positive pour $x>0$.
Son minimum vaut zéro : pour tout réel $x$,
\[
e^x\ge1+x,
\]
avec égalité seulement en zéro. Cette inégalité implique $e>1$.
\begin{center}\begin{tikzpicture}
\begin{axis}[width=10cm,height=5.5cm,axis lines=middle,xmin=-3,xmax=2,ymin=-1,ymax=7.5,xlabel={$x$},ylabel={$y$},legend pos=north west]
\addplot[accent,thick,domain=-3:2,samples=100]{exp(x)};\addlegendentry{$e^x$}
\addplot[black,dashed,domain=-2:2]{1+x};\addlegendentry{$1+x$}
\end{axis}\end{tikzpicture}\end{center}
Pour $x\to+\infty$, $e^x\ge1+x$ donne une croissance sans borne. Pour $x\to-\infty$, écrire $e^x=1/e^{-x}$ donne une limite nulle ; la valeur zéro n'est jamais atteinte.
\begin{finchapitre}
\section{Dériver et étudier}
La composition affine donne $[e^{ax+b}]'=ae^{ax+b}$. Pour $k>0$, $e^{kt}$ croît et $e^{-kt}$ décroît.
Leur valeur en zéro vaut $1$ ; le paramètre $k$ règle la vitesse d'évolution.
\begin{exemple}
Pour $f(x)=(x+1)e^{-x}$, $f'(x)=e^{-x}-(x+1)e^{-x}=-xe^{-x}$. Le signe de la dérivée est celui de $-x$. Le maximum global est $f(0)=1$, puisque la fonction croît avant zéro et décroît après.
\end{exemple}
\section{Liaison avec les suites}
Pour $u_n=e^{an+b}$, $u_{n+1}=e^a u_n$ : la suite est géométrique de raison positive $e^a$.
Une suite géométrique de raison négative alterne de signe et ne se prolonge donc pas par une exponentielle réelle positive de cette forme.


\section{Utilité et prolongements}
L'exponentielle décrit des évolutions dont la variation relative instantanée est constante : croissance idéale, décroissance radioactive ou relaxation d'un écart. Les hypothèses du modèle seront étudiées au chapitre suivant.

En licence, elle résout les équations différentielles linéaires ; pour des matrices, l'exponentielle définie par une série jouera un rôle analogue. En master, elle intervient dans les semi-groupes d'évolution, les transformées et les lois de probabilité. La loi $e^{x+y}=e^xe^y$ exprime déjà la composition de deux durées d'évolution.
\begin{retenir}
L'exponentielle est définie sur tous les réels, strictement positive et croissante. Ses lois de calcul résultent de son équation différentielle ; l'existence est construite plus loin, sans raisonnement circulaire.
\end{retenir}
\end{finchapitre}
\end{document}
